\section{Metodología}
\label{sec:metodologia}

\subsection{Marco metodológico: CRISP-DM}
\label{subsec:crispdm}

El proyecto se condujo con \emph{CRISP-DM} (\emph{Cross Industry Standard
Process for Data Mining}) como metodología principal \cite{wirth2000crisp}.
Las seis fases del ciclo se ejecutaron de forma iterativa: el proyecto
registró 134 ciclos experimentales y varias veces una fase posterior devolvió
el trabajo a una anterior. El caso más claro fue el hallazgo de la fuga por
pseudo-autores: apareció al auditar la validación y obligó a rehacer la
preparación de los datos (fusión de grupos) antes de continuar con el
modelado. La Tabla~\ref{tab:crispdm} resume el mapeo entre las fases, el
trabajo realizado y la evidencia que lo respalda.

\begin{table*}[t]
\centering
\caption{Las seis fases de CRISP-DM, el trabajo realizado y la evidencia
principal. El ciclo se recorrió varias veces; la tabla resume la ejecución
completa.}
\label{tab:crispdm}
\small
\begin{tabularx}{\textwidth}{@{}>{\raggedright\arraybackslash}p{2.1cm}X>{\raggedright\arraybackslash}p{3.1cm}@{}}
\toprule
Fase & Qué se hizo & Evidencia principal \\
\midrule
Comprensión del negocio & Se definió la tarea multietiqueta (HS, TR, AG), se registraron las reglas (sin transformers, sin etiquetas del test), se identificó la métrica con la sonda E00 (LB 0.51494) y se fijaron los criterios de éxito. & Sección I; \path{docs/competition_rules.md}; decisión D-009. \\
Comprensión de los datos & Se perfilaron 1\,731 tuits: prevalencias, perfiles conjuntos, longitudes, $\phi$, near-duplicados, pseudo-autores, asociación identificador--etiqueta y ambigüedad semántica de TR. & Sección III; \texttt{reports/eda/}; cuaderno, sección 2. \\
Preparación de los datos & Se fusionaron near-duplicados y pseudo-autores en grupos, se congelaron las particiones, se construyó V2\_triple y se verificó la ausencia de fuga en cada pliegue. & Sección IV; \path{features/}; decisiones D-004 y D-015. \\
Modelado & 134 ciclos experimentales: preprocesamiento, representación, modelos lineales, umbrales, ensembles, aprendizaje profundo desde cero, staging y pares compensados. & Secciones~\ref{sec:resultados} y~\ref{sec:arquitectura}; \texttt{experiments.csv}; \texttt{reports/}. \\
Evaluación & CCV agrupada en dos particiones, bootstrap, McNemar, análisis de error por rebanadas, holdout congelado y sondas del leaderboard. & Sección~\ref{sec:discusion}; \path{informe/datos/estadisticas.json}. \\
Despliegue & Reentrenamiento sobre el 100\,\% de \emph{train}, staging, envío, verificación SHA e idempotencia; publicación del cuaderno y de los artefactos. & Sección~\ref{sec:reproducibilidad}; \texttt{submissions/}; cuaderno. \\
\bottomrule
\end{tabularx}
\end{table*}

La gobernanza experimental complementó el ciclo. Cada experimento tuvo un
script propio y nuevo (una variación no edita el script anterior), un
pre-registro con pregunta, hipótesis y criterio de adopción, y una fila en
\texttt{experiments.csv}. La regla de adopción era exigente: una mejora local
debía superar $+0.005$ y sostenerse en las dos particiones; un candidato con
sonda debía mejorar además el leaderboard, y una mejora solo local no bastaba
(P109 es el contraejemplo). Cada corrida importante re-ejecutaba los
controles de determinismo (V03 exacto y P110 a $10^{-6}$) y el test nunca se
usó para entrenar ni para ajustar el modelo base. Esta disciplina explica por
qué el informe conserva los resultados negativos con el mismo detalle que los
positivos: son parte de la evidencia.

\subsection{Protocolo de validación agrupada}

La unidad de partición no es la fila sino el \emph{grupo fusionado}: la unión
de near-duplicados (similitud coseno de n-gramas de caracteres $\ge 0.85$) y
pseudo-autores (primera mención \texttt{@usuario}). Sobre los 1\,731 tuits se
construyen 1\,645 grupos. Con ellos se separa un desarrollo de 1\,483 filas y
un holdout congelado de 248 mediante \texttt{StratifiedGroupKFold} de siete
pliegues (semilla 42), estratificando por la combinación HS$\cdot$TR$\cdot$AG.
El desarrollo se evalúa con dos particiones de cinco pliegues: la canónica
(semilla 42) y la partición 47 (semilla 43), que actúa como réplica
independiente. Todas las métricas locales de este artículo son
\emph{out-of-fold} (OOF) sobre las 1\,483 filas de desarrollo. La línea base
del proyecto, V03, es la unión TF-IDF \texttt{word(1,2)}+\texttt{char\_wb(3,5)}
con regresión logística \textsc{liblinear} ($C=2.0$), umbrales 0.44/0.51/0.46 y
sin regla ascendente; su $F_1$ OOF es 0.777211 y su LB 0.78306. El holdout de
248 filas se reservó como guarda congelada y no se usó para las decisiones de
la serie final; por eso las cifras que se reportan son OOF sobre desarrollo.

El protocolo impone tres controles. Primero, los vectorizadores se ajustan
solo con el pliegue de entrenamiento y se aplican al de validación; el test
nunca participa. Segundo, se verifica en cada pliegue que ningún grupo
fusionado cruce la frontera entrenamiento/validación. Tercero, cada corrida
re-ejecuta dos controles de determinismo: la reconstrucción de la línea base
V03 debe coincidir exactamente (diferencia 0.0) y la de P110 dentro de
$10^{-6}$. Estos controles detectan deriva de entorno antes de comparar
candidatos \cite{kohavi1995study,kaufman2012leakage}.

\subsection{Representación subléxica multiescala}

La representación campeona, que llamamos V2\_triple, concatena tres bloques
TF-IDF \cite{joachims1998text}: palabras con unigramas y bigramas
(\texttt{word(1,2)}), n-gramas de caracteres con límites de palabra
(\texttt{char\_wb(3,5)}) y n-gramas de caracteres libres
(\texttt{char(1,6)}). Para cada bloque, el peso de un término $t$ en un
documento $d$ es
\begin{equation}
w_{t,d} = \bigl(1+\log \mathrm{tf}_{t,d}\bigr)
\left(\log\frac{1+N}{1+\mathrm{df}_t}+1\right),
\label{eq:tfidf}
\end{equation}
donde $\mathrm{tf}_{t,d}$ es la frecuencia del término en el documento (con
$w_{t,d}=0$ si el término no aparece), seguido de normalización $L_2$; se usa \texttt{min\_df}=2 y
\texttt{sublinear\_tf}=True. La unión produce 86\,275 características sobre el
corpus completo. El bloque de palabras aporta la señal semántica; los bloques
de caracteres aportan robustez frente a errores ortográficos, alargamientos
vocálicos, elisiones y jerga, que son precisamente los fenómenos que un
vocabulario de palabras no puede representar \cite{cavnar1994ngram,mcnamee2004character}.

\subsection{Modelos, ensamble y calibración}

Para cada etiqueta se ajusta una regresión logística con pesos de clase
balanceados,
\begin{equation}
p_j(x)=\sigma\bigl(\beta_{j0}+\beta_j^{\top}x\bigr), \qquad
\sigma(z)=\frac{1}{1+e^{-z}},
\label{eq:logit}
\end{equation}
con regularización $C$ adaptada por etiqueta (2.0 para HS, 1.8 para TR y 1.5
para AG). El solver es \textsc{liblinear} \cite{fan2008liblinear}. Sobre la
misma matriz se ajustan dos modelos adicionales: un \texttt{SGDClassifier}
con pérdida logarítmica y $\alpha=5\times10^{-6}$ \cite{bottou2010large}, y
una regresión logística SAGA con penalización ElasticNet
($\ell_1$-ratio 0.05) \cite{defazio2014saga,zou2005regularization}. El
ensamble, que llamamos Tri-Optimizer, combina las probabilidades
\begin{equation}
p = 0.80\,p^{\mathrm{LR}} + 0.05\,p^{\mathrm{SGD}} + 0.15\,p^{\mathrm{SAGA}},
\label{eq:ensemble}
\end{equation}
y las calibra con un escalado de temperatura en el espacio log-odds
\cite{guo2017calibration}:
\begin{equation}
\tilde p_j = \sigma\!\left(\frac{\operatorname{logit}(p_j)}{T}\right),
\qquad T=1.04.
\label{eq:temperatura}
\end{equation}
La temperatura se fijó en la fase de desarrollo y se congeló; no se ajustó
con el test. La combinación de tres optimizadores con sesgos inductivos
distintos busca reducir la varianza sin renunciar a la interpretabilidad de
un modelo lineal.

\subsection{Decisión: umbrales, regla ontológica y pares compensados}

La binarización usa un umbral por etiqueta,
\begin{equation}
\hat y_j = \mathbb{1}\left[\tilde p_j \ge \tau_j\right], \qquad
\tau = (0.44,\;0.51,\;0.46),
\label{eq:umbral}
\end{equation}
y después aplica la regla ascendente de HatEval: si $TR=1$ y $HS=0$, se
fuerza $HS=1$. Los umbrales se fijaron en la fase de calibración (desarrollo y
leaderboard público) con un criterio de $F_1$ por columna
\cite{lipton2014thresholding,nan2012optimizing} y se congelaron antes de la
serie final; no se usaron etiquetas del test. La métrica oficial es
\begin{equation}
\begin{split}
M &= \frac{1}{3}\sum_{j\in\{\mathrm{HS},\mathrm{TR},\mathrm{AG}\}} F_1^{(j)},\\
F_1^{(j)} &= \frac{2TP_j}{2TP_j+FP_j+FN_j}.
\end{split}
\label{eq:f1}
\end{equation}

Sobre esa base actúa la capa de decisión de la serie P120--P133. Su
componente principal son los \emph{pares compensados}: un lote de reglas
léxicas que rescata verdaderos positivos y suprime falsos positivos, con la
restricción de que el número de celdas activadas y desactivadas sea el mismo
en cada columna. La propiedad central se resume en la siguiente proposición.

\begin{proposicion}[Economía del par compensado]
\label{prop:par}
Sea una columna con $TP$, $FP$ y $FN$ celdas, y sea $D=2TP+FP+FN$. Si un
lote compensado rescata $a$ verdaderos positivos ($FN\to TP$) y suprime $a$
falsos positivos ($FP\to TN$) en esa columna, entonces
\begin{equation}
F_1' = F_1 + \frac{2a}{D} = F_1 + a\,\frac{F_1}{TP}.
\label{eq:par}
\end{equation}
El número de positivos de la columna no cambia.
\end{proposicion}

\begin{proof}
Tras el lote, $TP'=TP+a$, $FP'=FP-a$ y $FN'=FN-a$, de modo que
$D'=2(TP+a)+(FP-a)+(FN-a)=D$. Entonces
$F_1'=2(TP+a)/D=F_1+2a/D$. Como $F_1=2TP/D$, se tiene $2/D=F_1/TP$, y
sustituyendo se obtiene \eqref{eq:par}.
\end{proof}

La proposición explica por qué la ganancia de un par es predecible y por qué
la columna minoritaria es la más sensible: a igualdad de $a$, el incremento
$F_1/TP$ es mayor cuando $TP$ es pequeño, que es el caso de TR. También
describe el escenario óptimo: los $a$ rescates corresponden a verdaderos
positivos del conjunto oculto y las $a$ supresiones a falsos positivos reales.
Si una intervención rescata un caso que los anotadores consideraron negativo o
suprime un positivo genuino, la compensación se rompe y la ganancia se anula o
se vuelve negativa; P129 es la evidencia empírica de ese caso. En la implementación, el desplazamiento se aplica en log-odds
($\pm 2.0$ o $\pm 2.05$) para no saturar la probabilidad; la versión P129,
con un desplazamiento forzado de $+4.0$, ilustra el costo de romper la
compensación (Sección~\ref{sec:resultados}).

El lote P133, el campeón, contiene nueve pares: cuatro en HS, dos en TR y
tres en AG. Las firmas léxicas se diseñaron para tener cero coincidencias en
entrenamiento y una coincidencia exacta en test; por construcción, no
cambian la métrica OOF. La capa mantiene el \emph{ancla de prevalencia}
$\mathbf{n}=(514,330,421)$, es decir, 1\,265 positivos (36.54\,\%). El ancla no
es una estimación poblacional independiente: es la tasa de positivos que
maximizó el leaderboard en la fase de calibración y se usa como restricción
de invariancia. Esa distinción se discute en la Sección~\ref{sec:discusion}.

\subsection{Análisis estadístico}

Las comparaciones entre modelos usan bootstrap agrupado con 2\,000 réplicas y
semilla 20260925: se remuestrean grupos completos, no filas, para conservar
la dependencia \cite{efron1979bootstrap}. Los intervalos son percentiles al
95\,\%. Para los aciertos pareados se aplica la prueba exacta de McNemar
sobre las celdas discordantes \cite{mcnemar1947note,dietterich1998approximate}.
Las prevalencias se acompañan de intervalos de Wilson
\cite{wilson1927probable}; la asociación entre etiquetas se mide con $\phi$ y
$\chi^2$. La sensibilidad al umbral se resume en curvas $F_1$--$\tau$ por
etiqueta, y el error se desglosa por franjas de longitud (terciles en 70 y
118 caracteres). No se aplican correcciones por comparaciones múltiples
porque el número de hipótesis confirmatorias es pequeño y están
pre-registradas; los valores $p$ se reportan como evidencia, no como un
cribado masivo \cite{demsar2006statistical}.

La Tabla~\ref{tab:decisiones} sintetiza las decisiones de diseño y la
evidencia que las sostiene.

\begin{table*}[t]
\centering
\caption{Decisiones de diseño, alternativas consideradas y evidencia que las
sostiene. Las cifras locales son OOF sobre 1\,483 filas; las del leaderboard
son el registro público de Kaggle.}
\label{tab:decisiones}
\small
\begin{tabularx}{\textwidth}{@{}>{\raggedright\arraybackslash}p{2.0cm}>{\raggedright\arraybackslash}p{2.6cm}X>{\raggedright\arraybackslash}p{2.3cm}@{}}
\toprule
Decisión & Alternativa & Evidencia & Consecuencia \\
\midrule
Tres clasificadores binarios & Una clase con seis combinaciones & 0 casos $TR=1,HS=0$; 26 casos $AG=1,HS=0$; $\phi$ alto & Se adopta, con regla ascendente \\
Representación V2\_triple & Solo palabras o unión de dos bloques & $+0.00516$ local y $+0.00549$ LB frente a V03 & Se adopta \\
Tri-Optimizer & Mezcla 95/5 LR+SGD (P110) & $+0.00222$ local y $+0.00090$ LB; 4/5 y 3/5 pliegues & Se adopta \\
Umbrales por etiqueta & Umbral único 0.5 & $+0.00313$ LB en la línea base & Se congelan 0.44/0.51/0.46 \\
Validación agrupada & Validación por filas & 1\,645 grupos (86 filas menos); 0 fugas verificadas & Se adopta en todo el proyecto \\
Pares compensados & Staging masivo o unilateral & P128 $+0.00999$; P129 $-0.00438$ & Lotes pequeños, contables y verificados en LB \\
\bottomrule
\end{tabularx}
\end{table*}

La Figura~\ref{fig:arquitectura} resume la arquitectura del sistema.

\begin{figure*}[t]
\centering
\begin{tikzpicture}[
  node distance=2.6mm,
  stage/.style={draw, rounded corners, align=center, font=\footnotesize,
    inner sep=3pt, minimum height=9mm, text width=1.7cm},
  note/.style={draw=none, align=center, font=\scriptsize, text=black!65},
  arrow/.style={-{Latex[length=1.6mm]}, line width=0.5pt}
]
\node[stage] (corpus) {Corpus\\1\,731 / 1\,154};
\node[stage, right=of corpus] (tfidf) {TF-IDF\\V2\_triple\\86\,275};
\node[stage, right=of tfidf] (model) {Tri-Optimizer\\80/5/15\\$T=1.04$};
\node[stage, right=of model] (dec) {Umbrales\\0.44/0.51/0.46\\+ HatEval};
\node[stage, right=of dec] (pairs) {Pares\\compensados\\ancla};
\node[stage, right=of pairs] (out) {Envío\\+ SHA256};
\draw[arrow] (corpus) -- (tfidf);
\draw[arrow] (tfidf) -- (model);
\draw[arrow] (model) -- (dec);
\draw[arrow] (dec) -- (pairs);
\draw[arrow] (pairs) -- (out);
\node[note, below=3mm of tfidf] (val) {validación agrupada\\OOF (1\,483)};
\node[note, below=3mm of pairs] (lb) {sondas LB\\47 envíos};
\draw[arrow, dashed] (val) -- (tfidf);
\draw[arrow, dashed] (lb) -- (pairs);
\end{tikzpicture}
\caption{Arquitectura del sistema. La validación agrupada gobierna la
representación y el modelo; las sondas del leaderboard gobiernan únicamente
la capa de pares compensados. La separación es deliberada: son dos fuentes de
evidencia distintas.}
\label{fig:arquitectura}
\end{figure*}
